% Notation and field-discipline section. Should be short and
% self-contained — it's both a reading aid for the rest of the paper
% and a normative statement of the catalog's conventions.

A bilinear algorithm computing the product of an \(n \times m\)
matrix \(A\) by an \(m \times p\) matrix \(B\) consists of three
factor matrices \(U \in K^{nm \times r}\), \(V \in K^{mp \times
r}\), \(W \in K^{np \times r}\) over a field (or ring) \(K\), such
that for every \(A, B\):
\[
  (AB)_{i,j} \;=\; \sum_{k=1}^{r}
  W_{(i,j),k}\;\Big(\sum_{a,b} U_{(a,b),k}\, A_{a,b}\Big)
  \Big(\sum_{a',b'} V_{(a',b'),k}\, B_{a',b'}\Big).
\]
The integer \(r\) is the \emph{rank} (a.k.a.\ \emph{multiplication
count}) of the algorithm. The \emph{addition count} is the number of
nonzero coefficients in \((U, V, W)\) beyond the bare product
multiplications, with conventions following Probert--Fischer 1976.

\paragraph{Shape notation.} A bare \(\nmpshape{n}{m}{p}\) is not
sufficient -- the rank depends on the algebra. We adopt:
\begin{itemize}[itemsep=2pt]
  \item \(\nmpfield{K}{n}{m}{p}\): non-commutative matmul over the
    field \(K\). E.g.\ \(\nmpfield{\Reals}{7}{7}{7} = 250\).
  \item \(\nmpfieldc{K}{n}{m}{p}\) (superscript-\(c\)): commutative
    matmul over \(K\). E.g.\ \(\nmpfieldc{\Reals}{3}{3}{3} = 21\)
    (Rosowski).
\end{itemize}
The value after \(=\) is always the \emph{rank} --- the number of
scalar multiplications, which is the quantity that compounds into an
\(\omega\) bound under recursion. Addition counts are tracked
separately (Section~\ref{sec:architecture}); on disk the rank appears
as the \code{r\{rank\}} filename token (a cosmetic label --- the
authoritative rank is the JSON content; additions are recomputed from
the factor matrices, not encoded in the name).

\paragraph{The fields we track.}
\begin{description}[itemsep=2pt]
  \item[\Integers, \Rationals, \Reals] Characteristic-zero and
    non-commutative-friendly (they lift to recursive matmul). We keep
    them \emph{distinct}, first-class fields --- not a single cluster ---
    and use the inclusion \(\Integers\subset\Rationals\subset\Reals\)
    only to compute which fields a scheme satisfies. Ternary-integer
    schemes (\(\Integers\)T, coefficients in \(\{-1,0,1\}\)) are a
    first-class sub-class of \(\Integers\).
  \item[\Complex] A complex extension. Real-valid schemes work
    immediately; in addition, AlphaEvolve's complex-coefficient
    schemes live here. Its \(\nmpfield{\Complex}{4}{4}{4} = 48\) was
    subsequently rationalised by Dumas--Pernet--Sedoglavic
    \cite{dps2025}, so that particular rank no longer separates
    \Complex\ from \(\Rationals/\Reals\) --- a reminder that
    field-separation examples are date-stamped claims.
  \item[\Ftwo] Characteristic-2. AlphaTensor's
    \(\nmpfield{\Ftwo}{4}{4}{4} = 47\) lives here. Integer-valid
    schemes \emph{do} reduce mod 2 to give \Ftwo\ schemes, but our
    current widening sweep does not yet surface that exhaustively
    -- see Section \ref{sec:openquestions}.
  \item[\Fthree] Characteristic-3 (GF(3)) --- ``ternary modular'',
    not to be confused with the ternary-\emph{integer} ZT sub-class of
    \Integers. We track it so the discipline is genuinely
    field-agnostic, though coverage is currently thin: a dedicated
    \Fthree\ search is open (Section \ref{sec:openquestions}).
\end{description}

\paragraph{Commutative axis.} Independent of the field choice, an
algorithm may or may not assume that the entries of \(A\) and \(B\)
commute. Commutative-only schemes (Waksman, Rosowski, Makarov 1986)
do \emph{not} lift to recursive matmul over non-commutative rings:
they save bilinear products by exploiting the freedom \(ab = ba\)
that disappears as soon as \(A\) and \(B\) are themselves matrix
blocks. They remain valid for scalar matmul. Rosowski 2019 makes the
recursive side of this explicit: a commutative scheme \emph{can} be
applied recursively, but only while the base ring stays commutative
(the linear combinations of products are themselves bilinear, which is
what permits the recursion). His cubic bound is
\(\nmpfieldc{K}{n}{n}{n} \le n(n+1)(n+2)/2 \approx n^3/2\)
(Theorem~5, \(n\) odd) --- about half the \(n^3\) scalar products, and
below the best non-commutative rank at small sizes
(\(\nmpfieldc{\Reals}{3}{3}{3}=21\) versus Laderman's
\(\nmpfield{\Reals}{3}{3}{3}=23\)). The win is only this constant
factor: because the recursion never leaves the commutative ring it
does \emph{not} lift to matrix-block recursion, so unlike Strassen's
\(\nmpfield{K}{2}{2}{2}=7\) it does not lower \(\omega\) over
non-commutative rings. The catalog flags every
commutative scheme with \code{"commutative": true} in the JSON; the
comparison tables of Section \ref{sec:tables} carry a separate
commutative column for each field. Cross-contamination silently
produces wrong ``wins'' against historical non-commutative tables;
the canonical example is comparing
\(\nmpfieldc{\Reals}{3}{3}{3} = 21\) (Rosowski, commutative) to
\(\nmpfield{\Reals}{3}{3}{3} = 23\) (Laderman, non-commutative)
as if they were the same metric.

We require every rank claim in the paper to carry both a field tag
and a commutativity tag. The reader can recover the metric from the
shape.
